\subsection{MULTIMODULES}

Let A, B be two rings and consider on a set E two left module structures with the same additive law and whose ring of operators are respectively A and B; let \(\mathcal{E}\) be the endomorphism ring of the additive group E and for all \(\alpha \in A\) (resp. \(\beta \in B\) ) let \(h_{\alpha}\) (resp. \(h_{\beta}^{\prime}\) ) denote the element \(x \mapsto \alpha x\) (resp. \(x \mapsto \beta x\) ) of \(\mathcal{E}\) . Clearly the three following properties are equivalent: (a) \(h_{\alpha} \circ h_{\beta}^{\prime} = h_{\beta}^{\prime} \circ h_{\alpha}\) for all \(\alpha\) and \(\beta\) ; (b) the image of A under the homomorphism \(a \mapsto h_{\alpha}\) is contained in \(\operatorname{Hom}_{\mathbf{B}}(\mathbf{E}, \mathbf{E})\) ; (c) the image of B under the homomorphism \(\beta \mapsto h_{\beta}^{\prime}\) is contained in \(\operatorname{Hom}_{\mathbf{A}}(\mathbf{E}, \mathbf{E})\) . When the A-module (resp. B-module) structure in question is a right module structure, the ring A (resp. B) must be replaced in (b) (resp. (c)) by \(\mathbf{A}^{0}\) (resp. \(\mathbf{B}^{0}\) ). The above properties can be expressed by saying that the two (left or right) module structures defined on E are compatible.

DEFINITION 13. Let \((\mathbf{A}_{\lambda})_{\lambda \in \mathbf{L}}, (\mathbf{B}_{\mu})_{\mu \in \mathbf{M}}\) be two families of rings; an \(((\mathbf{A}_{\lambda}), (\mathbf{B}_{\mu}))\) -multimodule (or multimodule over the families of rings \((\mathbf{A}_{\lambda})_{\lambda \in \mathbf{L}}, (\mathbf{B}_{\mu})_{\mu \in \mathbf{M}})\) is a set E with, for each \(\lambda \in \mathbf{L}\) , a left \(\mathbf{A}_{\lambda}\) -module structure and, for each \(\mu \in \mathbf{M}\) , a right \(\mathbf{B}_{\mu}\) -module structure, all these module structures being compatible with one another.

When the family \((\mathbf{B}_{\mu})\) (resp. \((\mathbf{A}_{\lambda})\) ) is empty, E is called a left (resp. right) multimodule. When \(\operatorname{Card}(\mathbf{L}) + \operatorname{Card}(\mathbf{M}) = 2\) , we say ``bimodule'' instead of ``multimodule''; it is then often convenient to consider (as can always be done by replacing a ring of operators by its opposite, cf.~no. 1) a bimodule as having a left module structure with respect to a ring A and a right module structure with respect to a ring B, the permutability of the laws then being expressed by the relation

\begin{enumerate}
\def\labelenumi{(\arabic{enumi})}
\setcounter{enumi}{48}
\tightlist
\item
\end{enumerate}

\[
\alpha (x \beta) = (\alpha x) \beta \quad \text { for } x \in \mathbf {E}, \alpha \in \mathbf {A}, \beta \in \mathbf {B}.
\]

It is then also said that \(\mathbf{E}\) is an (A, B)-bimodule.

Two multimodule structures on a set E are said to be compatible if all the module structures on E which define one or the other of these multimodule structures are compatible with one another.

Examples. (1) On a ring A the module structures of \(A_{s}\) and \(A_{d}\) are compatible and A can therefore be considered canonically as an (A, A)-bimodule.

\begin{enumerate}
\def\labelenumi{(\arabic{enumi})}
\setcounter{enumi}{1}
\tightlist
\item
  A left A-module E has a canonical left module structure over the ring \(\operatorname{End}_{\mathbf{A}}(\mathbf{E})\) and the A-module and \(\operatorname{End}_{\mathbf{A}}(\mathbf{E})\) -module structures on E are compatible.
\end{enumerate}

Clearly when E is a multimodule over two families \((\mathbf{A}_{\lambda})_{\lambda\in L},(\mathbf{B}_{\mu})_{\mu\in M}\) of rings, E is also a multimodule over any two subfamilies \((\mathbf{A}_{\lambda})_{\lambda\in L},(\mathbf{B}_{\mu})_{\mu\in M'}\) , where the \(A_{\lambda}\) -module and \(B_{\mu}\) -module structures for \(\lambda\in L'\) and \(\mu\in M'\) are those initially given.

Since multimodules are particular examples of commutative groups with operators, the results of nos. 2 to 10 (cf.~no. 10, Remark) can be applied to them; in particular, if E, F are two \(((A_{\lambda}), (B_{\mu}))\) -multimodules, a homomorphism \(u: E \to F\) is a mapping which is an \(A_{\lambda}\) -homomorphism for all \(\lambda \in L\) and a \(B_{\mu}\) -homomorphism for all \(\mu \in M\) . The stable subgroups of an \(((A_{\lambda}), (B_{\mu}))\) -multimodule are \(((A_{\lambda}), (B_{\mu}))\) -multimodules (called submultimodules), as also are the quotients by such subgroups (called quotient multimodules); similarly for products and direct sums.

Let E be an \(((A_{\lambda}), (B_{\mu}))\) -multimodule and for each \(\lambda \in L\) (resp. \(\mu \in M\) ) let \(\phi_{\lambda}: A_{\lambda}' \to A_{\lambda}\) (resp. \(\psi_{\mu}: B_{\mu}' \to B_{\mu}\) ) be a ring homomorphism; clearly the \(A_{\lambda}'\) -module structures associated with the \(\phi_{\lambda}\) and the \(A_{\lambda}\) -module structures given on E and the \(B_{\mu}'\) -module structures associated with the \(\psi_{\mu}\) and the \(B_{\lambda}\) -module structures given on E are compatible with one another and hence define on E an \(((A_{\lambda}')\) , \((B_{\mu}')\) )-multimodule structure, said to be associated with the given \(((A_{\lambda})\) , \((B_{\mu})\) )-multimodule structure and the \(\phi_{\lambda}\) and \(\psi_{\mu}\) .

If E, F are two \(((A_{\lambda}), (B_{\mu}))\) -multimodules, the additive group of homomorphisms of E into F is denoted by \(\operatorname{Hom}_{(A_{\lambda}), (B_{\mu})}(E, F)\) (or simply \(\operatorname{Hom}(E, F)\) ). Formulae (6) to (8) of no. 2 are obviously valid for \(((A_{\lambda}), (B_{\mu}))\) -multimodule homomorphisms and, in particular, \(\operatorname{Hom}(E, E) = \operatorname{End}(E)\) has a ring structure;

moreover \(\operatorname{Hom}(\mathbf{E},\mathbf{F})\) has a canonical left \(\operatorname{End}(\mathbf{F})\) -module structure and right \(\operatorname{End}(\mathbf{E})\) -module structure, these two structures being compatible; in other words, \(\operatorname{Hom}(\mathbf{E},\mathbf{F})\) has a canonical (End(F), End(E))-bimodule structure.

Suppose now that E has a multimodule structure whose rings of left (resp. right) operators are on the one hand the \(A_{\lambda}\) for \(\lambda \in L\) (resp. the \(B_{\mu}\) for \(\mu \in M\) ) and on the other the rings of another family \((\mathbf{A}_{\lambda'}')_{\lambda' \in L'}\) (resp. \((\mathbf{B}_{\mu'}')_{\mu' \in M'}\) ). Suppose similarly that F has a multimodule structure whose rings of left (resp. right) operators are on the one hand the \(A_{\lambda}\) for \(\lambda \in L\) (resp. the \(B_{\mu}\) for \(\mu \in M\) ) and on the other the rings of another family \((\mathbf{A}_{\lambda''})_{\lambda'' \in L''}\) (resp. \((\mathbf{B}_{\mu''})_{\mu'' \in M''}\) ); to abbreviate we shall say that E is an \(((\mathbf{A}_{\lambda}), (\mathbf{A}_{\lambda'}'); (\mathbf{B}_{\mu}), (\mathbf{B}_{\mu'})\) -multimodule and F an \(((\mathbf{A}_{\lambda}), (\mathbf{A}_{\lambda''}); (\mathbf{B}_{\mu}), (\mathbf{B}_{\mu''}))\) -multimodule. Consider E and F as \(((\mathbf{A}_{\lambda}), (\mathbf{B}_{\mu}))\) -multimodules, thus restricting the operators to the subfamilies \((\mathbf{A}_{\lambda})\) and \((\mathbf{B}_{\mu})\) . By what was said at the beginnings of this no., the multimodule structures given on E and F define canonically ring homomorphisms

\[
\mathrm{A} _ {\lambda^ {\prime}} ^ {\prime} \rightarrow \operatorname{End} _ {(\mathrm{A} _ {\lambda}), (\mathrm{B} _ {\mu})} (\mathrm{E}), \quad \mathrm{B} _ {\mu^ {\prime}} ^ {\prime 0} \rightarrow \operatorname{End} _ {(\mathrm{A} _ {\lambda}), (\mathrm{B} _ {\mu})} (\mathrm{E}),
\]

\[
\mathrm{A} _ {\lambda^ {\prime \prime}} ^ {\prime \prime} \rightarrow \operatorname{End} _ {(\mathrm{A} _ {\lambda}), (\mathrm{B} _ {\mu})} (\mathrm{F}), \quad \mathrm{B} _ {\mu^ {\prime \prime}} ^ {\prime \prime} \rightarrow \operatorname{End} _ {(\mathrm{A} _ {\lambda}), (\mathrm{B} _ {\mu})} (\mathrm{F});
\]

moreover, two elements of \(\operatorname{End}_{(\mathbf{A}_{\lambda}), (\mathbf{B}_{\mu})}(\mathbf{E})\) (resp. \(\operatorname{End}_{(\mathbf{A}_{\lambda}), (\mathbf{B}_{\mu})}(\mathbf{F})\) ) respective images of elements of two distinct rings among the \(\mathbf{A}_{\lambda'}'\) or the \(\mathbf{B}_{\mu'}^{0}\) (resp. the \(\mathbf{A}_{\lambda''}\) or the \(\mathbf{B}_{\mu''}^{0}\) ) are permutable; it follows that the above homomorphisms define on \(\operatorname{Hom}_{(\mathbf{A}_{\lambda}), (\mathbf{B}_{\mu})}(\mathbf{E}, \mathbf{F})\) a multimodule structure whose rings of left operators are the \(\mathbf{A}_{\lambda''}\) ( \(\lambda'' \in \mathbf{L}''\) ) and the \(\mathbf{B}_{\mu'}'\) ( \(\mu' \in \mathbf{M}'\) ) and whose rings of right operators are the \(\mathbf{A}_{\lambda'}'\) ( \(\lambda' \in \mathbf{L}'\) ) and the \(\mathbf{B}_{\mu''}^{0}(\mu'' \in \mathbf{M}'')\) .

If now \(\mathbf{E}'\) is an \(((\mathbf{A}_{\lambda}), (\mathbf{A}_{\lambda'}^{\prime}) ; (\mathbf{B}_{\mu}), (\mathbf{B}_{\mu'}^{\prime}))\) -multimodule and \(\mathbf{F}'\) an \(((\mathbf{A}_{\lambda}), (\mathbf{A}_{\lambda''}^{\prime}) ; (\mathbf{B}_{\mu}), (\mathbf{B}_{\mu''}^{\prime}))\) -multimodule, \(\mathrm{Hom}_{(\mathbf{A}_{\lambda}), (\mathbf{B}_{\mu})}(\mathbf{E}', \mathbf{F}')\) is an

\[
\left(\left(\mathrm{A} _ {\lambda^ {\prime \prime}} ^ {\prime \prime}\right), \left(\mathrm{B} _ {\mu^ {\prime}} ^ {\prime}\right): \left(\mathrm{A} _ {\lambda^ {\prime}} ^ {\prime}\right), \left(\mathrm{B} _ {\mu^ {\prime \prime}} ^ {\prime \prime}\right)\right) - \text { multimodule };
\]

if \(u: \mathbf{E}' \to \mathbf{E}\) , \(v: \mathbf{F} \to \mathbf{F}'\) are multimodule homomorphisms,

\[
\operatorname{Hom} (u, v): \operatorname{Hom} _ {\left(\mathrm{A} _ {\lambda}\right), \left(\mathrm{B} _ {\mu}\right)} (\mathrm{E}, \mathrm{F}) \rightarrow \operatorname{Hom} _ {\left(\mathrm{A} _ {\lambda}\right), \left(\mathrm{B} _ {\mu}\right)} \left(\mathrm{E} ^ {\prime}, \mathrm{F} ^ {\prime}\right)
\]

is defined as in no. 2 and is a multimodule homomorphism.

Remarks. (1) Let F be an A-module and C the centre of the ring A; as central homotheties commute with all homotheties, F has a bimodule structure whose rings of left operators are A and C. If E is another A-module, \(\operatorname{Hom}_{\mathbf{A}}(\mathbf{E}, \mathbf{F})\) has therefore a canonical C-module structure (where, for \(f \in \operatorname{Hom}_{\mathbf{A}}(\mathbf{E}, \mathbf{F})\) and \(\gamma \in C\) , \(\gamma f\) is the homomorphism \(x \mapsto \gamma f(x)\) ); if \(E'\) , \(F'\) are two A-modules, \(u: E' \to E\) , \(v: F \to F'\) two A-homomorphisms, the mapping \(\operatorname{Hom}(u, v)\) is C-linear.

\begin{enumerate}
\def\labelenumi{(\arabic{enumi})}
\setcounter{enumi}{1}
\tightlist
\item
  Let E be a left A-module; as A has a canonical (A, A)-bimodule structure, so has the direct sum \(\mathbf{A}^{(\mathrm{T})}\) for any indexing set T; by the above, \(\mathrm{Hom}_{\mathbf{A}}(\mathbf{A}_{s}^{\mathrm{(T)}},\mathbf{E})\) has a canonical left A-module structure arising from the right A-module structure on \(\mathbf{A}_{s}^{\mathrm{(T)}}\) : for \(f\in\mathrm{Hom}_{\mathbf{A}}(\mathbf{A}_{s}^{\mathrm{(T)}},\mathbf{E})\) and \(\alpha\in\mathbf{A},\alpha f\) is the linear mapping \(x\mapsto f(x\alpha)\) . Corollary 2 to Proposition 17 of no. 11 defines a canonical mapping \(j_{\mathbf{E},\mathbf{T}}\) from the product module \(\mathbf{E}^{\mathrm{T}}\) to \(\mathrm{Hom}_{\mathbf{A}}(\mathbf{A}_{s}^{\mathrm{(T)}},\mathbf{E})\) , the image under \(j_{\mathbf{E},\mathbf{T}}\) of a family \((x_{t})_{t\in\mathbf{T}}\) being the linear mapping \(f:\mathbf{A}_{s}^{\mathrm{(T)}}\to\mathbf{E}\) such that \(f(e_{t})=x_{t}\) for all \(t\in\mathbf{T}\) (where \((e_{t})\) is the canonical basis of \(\mathbf{A}_{s}^{\mathrm{(T)}}\) ); it is known (loc. cit.) that \(j_{\mathbf{E},\mathbf{T}}\) is bijective and it follows from the definition given above of the A-module structure on \(\mathrm{Hom}_{\mathbf{A}}(\mathbf{A}_{s}^{\mathrm{(T)}},\mathbf{E})\) that \(j_{\mathbf{E},\mathbf{T}}\) is A-linear. Finally, if \(u:\mathbf{E}\to\mathbf{F}\) is an A-module homomorphism, the diagram
\end{enumerate}

\[
\begin{array}{c} \mathrm {E^ {T}} \xrightarrow {j _ {\mathrm{E,T}}} \operatorname{Hom} _ {\mathrm{A}} (\mathrm{A} _ {s} ^ {(\mathrm{T})}, \mathrm{E}) \\ u ^ {\mathrm{T}} \Biggl \downarrow \qquad \qquad \qquad \qquad \qquad \qquad \qquad \qquad \qquad \qquad \qquad \qquad \qquad \qquad \qquad \qquad \qquad \qquad \qquad \qquad \qquad \qquad \qquad \qquad \qquad \qquad \qquad \text {Hom(1,u)} \\ \mathrm {F^ {T}} \xrightarrow {j _ {\mathrm{F,T}}} \operatorname{Hom} _ {\mathrm{A}} (\mathrm{A} _ {s} ^ {(\mathrm{T})}, \mathrm{F}) \end{array}\tag{50}
\]

is commutative.

Note that when T consists of a single element, \(j_{E}:E \to \operatorname{Hom}_{A}(A_{s}, E)\) is just the mapping \(x \mapsto \theta_{x}\) defined in no. 2, Example 1.

